\documentclass[10pt,a4paper]{amsart}
\usepackage[margin=19mm]{geometry}
\usepackage{amsmath,amssymb,mathtools}
\usepackage[scr=rsfs,cal=euler]{mathalfa}
\usepackage{xfrac}
\usepackage[unicode,hidelinks,bookmarks=false]{hyperref}
\newcommand{\ie}{i.e.\ }
\newcommand{\cf}{cf.\ }
\newcommand{\resp}{respectively\ }
\newcommand{\Subsection}{\subsection}
\providecommand{\nobreakdash}{\nobreak\hbox{-}}
\setlength{\emergencystretch}{2em}
\multlinegap0pt
\usepackage{bm}
\usepackage{bbm}
\usepackage{dsfont}
\usepackage{xspace}
\usepackage{cancel}
\usepackage{mathtools}
\usepackage{amsthm}
\makeatletter
\@ifundefined{thm}{\newtheorem{thm}{Thm}}{}
\@ifundefined{corollary}{\newtheorem{corollary}[thm]{Corollary}}{}
\@ifundefined{theorem}{\newtheorem{theorem}[thm]{Theorem}}{}
\makeatother
\newcommand\Sol{{\mathcal {S}}\!{\text{\it{ol}}}\,}
\title[A generating operator for Rankin-Cohen brackets]{A generating operator for Rankin-Cohen brackets}
\author{Toshiyuki Kobayashi, Michael Pevzner}
\date{Source: arXiv:2306.16800v2 (CC BY 4.0)}
\begin{document}
\maketitle

\noindent\emph{Source and licence.} A generating operator for Rankin-Cohen brackets. Version arXiv:2306.16800v2 (CC BY 4.0), distributed under the Creative Commons Attribution 4.0 International licence, as recorded in the arXiv licence metadata for this exact version and shown on the abstract page https://arxiv.org/abs/2306.16800v2. Full rights notice: licenses/arxiv-2306.16800-CC-BY-4.0.txt.\par\smallskip

\noindent\textit{Editorial scope.} Counted unit: the Theorem whose source label is \texttt{thm:Tlsol} in the article (source0.tex lines 1266--1289, source label \texttt{thm:Tlsol}), delivered together with the article's own complete original proof of it (source0.tex lines 1303--1403, one \texttt{proof} environment of 101 lines). No mathematics is written, repaired or re-derived: statement and proof are the article's own text line for line. The proof is detached from the statement in the source: the article states the result at lines 1266--1289 and prints its proof at lines 1303--1403, and the proof environment's own header names the statement, so the association is the article's own. Required context retained uncounted: eqn:QTaylor (lines 553--558); eqn:23031847 (lines 638--646); cor:Luminy0316 (lines 763--774). Every definition, display and label of those lines is retained unchanged. Omitted from this package and not redistributed: 1--552, 559--637, 647--762, 775--1265, 1290--1302, 1404--2385. Nothing inside a retained excerpt is omitted or altered: every retained body is delivered exactly as the article's lines are written, so no omission receipt of any kind accompanies this package. Citations: the retained excerpts carry no citation command at all, so no bibliography entry of the article is transcribed and no reference is redistributed. Reference adaptation: none was needed and none was made; every reference inside the delivered unit resolves inside it, so no unresolved reference or citation is produced. Printed slips kept literal: no printed word, symbol, hypothesis or number of the counted statement or of its proof was altered. Layout and class: the article's own class is \texttt{amsart} with packages including amsmath, amssymb, amsfonts, amscd, pifont, amsthm, amsxtra, latexsym, xy, here, graphicx, psfrag; the wrapper is the lane's pinned \texttt{amsart} editorial wrapper at 10pt on A4, which restates the theorem-like environments the retained text needs and adds the article's own macro definitions that the retained text needs (\texttt{\textbackslash Sol}), with extra wrapper packages (bm, bbm, dsfont, xspace, cancel, mathtools). The delivered numbering is the unit's own. Assets: no retained span contains a figure environment, a graphics-inclusion command, a PDF-page inclusion, a table or a TikZ picture; no image file is copied into this package, the package declares no asset dependency and no image file is redistributed with it. Licence: version arXiv:2306.16800v2 of this article is distributed under the Creative Commons Attribution 4.0 International licence, as recorded in the arXiv licence metadata for this exact version and shown on the abstract page https://arxiv.org/abs/2306.16800v2. A full rights notice is delivered at licenses/arxiv-2306.16800-CC-BY-4.0.txt. Characters: every retained excerpt is pure ASCII, so the delivered engine, XeLaTeX with its default Latin Modern text font, reports no missing character.
\begin{equation}
\label{eqn:QTaylor}
  \frac 1 Q
  =
  \sum_{\ell=0}^\infty \frac{(-1)^{\ell}(\zeta_1-\zeta_2)^{\ell}t^{\ell}}{(\zeta_1-z)^{\ell+1}(\zeta_2-z)^{\ell+1}}.  
\end{equation}

\begin{equation}
\label{eqn:23031847}
P:=
(\zeta_1-\zeta_2)^2 
\frac{\partial^2}{\partial \zeta_1 \partial \zeta_2}
-
(\zeta_1-\zeta_2) 
(\frac{\partial}{\partial \zeta_1}-\frac{\partial}{\partial \zeta_2}).  
\end{equation}

\begin{corollary}
\label{cor:Luminy0316}
Let $\ell \in {\mathbb{N}}$.  
Then the following two conditions on $f \in {\mathcal{O}}(D \times D)$ are equivalent:
\par\noindent
{\rm{(i)}}\enspace
$f \in \Sol(D \times D)_{\ell}$, 
\par\noindent
{\rm{(ii)}}\enspace
$T f (z,t)$ is of the form $t^{\ell} \varphi (z)$
 for some $\varphi \in {\mathcal{O}}(D)$.    
\end{corollary}

\begin{theorem}
\label{thm:Tlsol}
Let $\ell \in {\mathbb{N}}$.  
\par\noindent
{\rm{(1)}}\enspace
$t^{-\ell}T \colon \Sol(D \times D)_{\ell} \overset \sim \rightarrow {\mathcal{O}}(D)$
 is a bijection.  
\par\noindent
{\rm{(2)}}\enspace
The inverse of $t^{-\ell}T$ is given 
 by the integral operator $\Psi_{\ell}$, 
 namely, 
$
   \Psi_{\ell} \colon {\mathcal{O}}(D)
   \overset \sim \rightarrow
   \Sol(D \times D)_{\ell}
$
 is a bijection
 and 
$
  t^{-\ell} T \circ \Psi_{\ell}
  = \frac{2^{2\ell+1}}{2 \ell+1} \operatorname{id}.
$
\end{theorem}

\begin{proof}
[Proof of Theorem \ref{thm:Tlsol}]
First, 
 we prove
$
  \operatorname{Image}\Psi_{\ell} \subset \Sol(D \times D)_{\ell}.  
$
Recall from \eqref{eqn:23031847} the definition of $P.$  
A direct computation shows 
\begin{align*}
&P(\Psi_{\ell}g)
+ \ell (\ell+1) \Psi_{\ell}g 
\\
=&
 -\frac 1 2 (\zeta_1 - \zeta_2)^{\ell+1} \int_{-1}^1 \frac{\partial}{\partial v}
 (g'(\frac{(\zeta_2- \zeta_1)v+(\zeta_1 + \zeta_2)}2)(1-v^2)^{\ell+1}) d v, 
\end{align*}
which vanishes for any $g \in {\mathcal{O}}(D)$.  
Hence $\Psi_{\ell}g$ is an eigenfunction of $P$
 for the eigenvalue $-\ell(\ell+1)$.  



Second, 
 we prove
 that the \lq\lq{generating operator}\rq\rq\ $T$ gives
 the inverse of $\Psi_{\ell}$
 up to scalar multiplication, 
 that is, 
\begin{equation}
\label{eqn:Tpsi}
   T(\Psi_{\ell}g)(z,t)=\frac{2^{2 \ell+1}}{2 \ell +1} t^{\ell}g(z)
\qquad\text{for any $g \in {\mathcal{O}}(D)$.  }
\end{equation}



To see \eqref{eqn:Tpsi}, 
 we observe from Corollary \ref{cor:Luminy0316}
 that $t^{-\ell}(T \Psi_{\ell}g)(z,t)$ does not depend 
 on the variable $t$
 because $\Psi_{\ell} g \in \Sol(D \times D)_{\ell}$.  
On the other hand, 
 it follows from the expansion \eqref{eqn:QTaylor}
 that the coefficient of $t^{\ell}$ in $(T f)(z,t)$
 is given by 
\[
  \frac{1}{(2 \pi \sqrt{-1})^2}
  \oint \oint \frac{(-1)^{\ell}(\zeta_1-\zeta_2)^{\ell} f(\zeta_1, \zeta_2)}{
(\zeta_1-z)^{\ell+1}(\zeta_2-z)^{\ell+1}} d \zeta_1 d \zeta_2.  
\]



Applying this to $f=\Psi_{\ell}g$, 
 one sees that $t^{-\ell}(T \Psi_{\ell} g)(z,t)$ is equal to 
\begin{align*}
  &\frac{(-1)^{\ell}}{(2 \pi \sqrt{-1})^2}
  \oint \oint 
  \frac
  {(\zeta_1-\zeta_2)^{2\ell}\int_{-1}^1 g(\frac{(\zeta_2- \zeta_1)v+(\zeta_1 + \zeta_2)}2)(1-v^2)^{\ell} d v}
  {(\zeta_1-z)^{\ell+1}(\zeta_2-z)^{\ell+1}} d \zeta_1 d \zeta_2
\\
=& \frac{(-1)^{\ell}}{(\ell !)^2}
   \left. \frac{\partial^{2\ell}}
               {\partial \zeta_1^{\ell} \partial \zeta_2^{\ell}}\right|_{\zeta_1=\zeta_2=z}
   ((\zeta_1-\zeta_2)^{2\ell}
    \int_{-1}^1 g(\frac{(\zeta_2- \zeta_1)v+(\zeta_1 + \zeta_2)}2)(1-v^2)^{\ell} d v).  
\end{align*}



An iterated use of the Leibniz rule develops the right-hand side
 as a sum of various derivatives, 
 among which the only non-vanishing term is 
\[ g(z)\frac{(2 \ell)!}{(\ell!)^2}
   \int_{-1}^1 (1-v^2)^{\ell} d v 
= \frac{2^{2\ell+1}}{2 \ell+1}g(z).  
\]
Thus we have shown \eqref{eqn:Tpsi}, 
 hence the injective morphism
 $t^{-\ell}T \colon \Sol(D \times D)_{\ell} \to {\mathcal{O}}(D)$
 is also surjective.  



Finally, 
 let us show the surjectivity of $\Psi_{\ell}$.  
For any $f \in \Sol(D \times D)_{\ell}$, 
 there exists $g \in {\mathcal{O}}(D)$
 such that
$(T f)(z,t)=t^{\ell}g (z)$
 by Corollary \ref{cor:Luminy0316}.  
Since the right-hand side equals 
 $\frac{2\ell+1}{2^{2\ell+1}} T \Psi_{\ell}(g)$
  by \eqref{eqn:Tpsi}, 
 one has 
 $f=\frac{2\ell+1}{2^{2\ell+1}} \Psi_{\ell}g$
 because $T$ is injective.  
Thus the surjectivity of $\Psi_{\ell}$ is shown.  
\end{proof}

\end{document}
